%%%%%%%%%%%%%%%%%%%%%%%%%%%%%%%%%%%%%%%%%%%%%%%%%%%%%%%%%%%%%%%%%%%%%%
% Lớp 12 | Các số đặc trưng ... ghép nhóm | Khoảng biến thiên, khoảng tứ phân vị của mẫu số liệu ghép nhóm
% Dạng 2D3?1-4: Câu hỏi tổng hợp
% Số câu: 9  (sinh tự động bởi tach_cau.py; từ nay thầy sửa trực tiếp file này)
%%%%%%%%%%%%%%%%%%%%%%%%%%%%%%%%%%%%%%%%%%%%%%%%%%%%%%%%%%%%%%%%%%%%%%

%%%%%%%%%%%%%%%% ĐÚNG/SAI %%%%%%%%%%%%%%%%

% nguồn: Lop_12_Dai_so_Chuong_3
\begin{ex}%[2D3H1-4]%[ma:7792c71e]
	Tìm hiểu thời gian sử dụng điện thoại trong tuần đầu tháng 6/2024 của kỳ nghỉ hè lớp chủ nhiệm.  GVCN thu được kết quả sau
	\begin{center}
		\begin{tabular}{|c|c|c|c|c|c|c|}
			\hline 
			Thời gian (giờ) & $\left[0;5\right)$ & $\left[5;10\right)$ & $\left[10;15\right)$ & $\left[15;20\right)$ & $\left[20;25\right)$ & $\left[25;30\right]$ \\ 
			 \hline
			Số học sinh & $2$ & $6$ & $8$ & $9$ &  $3$ & $2$ \\
			\hline 
		\end{tabular}
	\end{center}
	 Khi đó
	 \choiceTF
		{\True Khoảng biến thiên của mẫu số liệu ghép nhóm này là $30$} 
		{\True Nhóm chứa tứ phân vị thứ  $3$ là $\left[15;20\right)$}
		{Số  trung bình của mẫu số liệu là $10$}
		{\True Khoảng  tứ phân vị của mẫu số liệu ghép nhóm này lớn hơn $10$}
	\loigiai{\begin{center}
	\begin{tabular}{|c|c|c|c|c|c|c|}
		 \hline
		Thời gian (giờ) & $\left[0;5\right)$ & $\left[5;10\right)$ &  $\left[10;15\right)$ & $\left[15;20\right)$ & $\left[20;25\right)$ & $\left[25;30\right]$ \\
		\hline 
		Số học sinh & $2$ & $6$ & $8$ & $9$ &  $3$ & $2$ \\
		\hline 
		Giá trị đại  diện & $2{,}5$ & $7{,}5$ & $12{,}5$ & $17{,}5$ & $22{,}5$ & $27{,}5$ \\
		 \hline
	\end{tabular}
\end{center}
\begin{itemchoice}
	\itemch Đúng.\\ 
	Khoảng biến thiên của mẫu số liệu ghép nhóm  là $R=30-0=30$.
	\itemch  Đúng.\\
	Vì tứ phân vị nhóm thứ 3 của mẫu  số liệu gốc là $Q_3=x_{23} \in \left[15;20\right)$.
	 \itemch Sai.\\
	Thời gian sử  dụng điện thoại trung bình của học sinh là\\
	$$\overline{x}=\dfrac{2\cdot 2{,}5+6\cdot 7{,}5+8\cdot 12{,}5+9\cdot  17{,}5+3\cdot 22{,}5+2\cdot 27{,}5}{30}=\dfrac{43}{3} \approx 14{,}3.$$
	 \itemch Đúng.\\
	Mẫu số liệu gốc  $Q_1=x_8\in \left[5;18\right),Q_3=x_{23} \in \left[15;20\right)$ từ đó ta có\\ $Q_1=5+\dfrac{\dfrac{30}{4}-2}{6}\cdot 5=9{,}58$; $Q_3=15+\dfrac{\dfrac{90}{4}-16}{9}\cdot 5=\dfrac{335}{18} \approx 18{,}61$.\\
	$\Rightarrow \triangle_Q=Q_3-Q_1=\dfrac{325}{36} \approx 9{,}03<  10$.
\end{itemchoice}
 }
\end{ex}

% nguồn: Lop_12_Dai_so_Chuong_3
\begin{ex}%[2D3H1-4]%[ma:a2613d41]
	 Hai thùng hàng $A$, $B$ đều chứa $25$ quả táo. Kết quả kiểm tra cân nặng của $25$ quả táo ở mỗi thùng $A$ và $B$ được cho ở bảng sau:
	\begin{center}
		\resizebox{0.8\linewidth}{!}{\begin{tabular}{|c|c|c|c|c|c|}
		 \hline
		Cân nặng (g)  & $[250;260)$ & $[260;270)$ & $[270;280)$ & $[280;290)$ & $[290;300)$\\
		 \hline
		Số  táo ở thùng $A$ & $2$ & $4$ & $12$ & $4$ & $3$\\
		 \hline
		Số táo ở thùng $B$ & $1$ & $3$ & $7$  & $10$ & $4$\\
		\hline 
	\end{tabular}}
\end{center}
 \choiceTF
	{\True Số trung bình của cân nặng các quả táo ở thùng $A$ nhỏ hơn số trung bình của cân nặng các quả táo ở thùng $B$} 
	{\True Lấy  ngẫu nhiên một quả táo từ thùng $A$. Xác suất để quả táo đó có cân nặng từ $280$ g trở lên là~$0{,}28$}
	{\True Lấy ngẫu nhiên một quả táo từ thùng $A$ và một quả táo từ thùng $B$. Xác suất để hai quả táo lấy ra đều nặng từ  $270$ g trở lên là~$0{,}6384$}
	{\True Số táo có cân nặng từ $280$ g trở lên ở thùng $B$ nhiều hơn số táo có cân nặng từ $280$  g trở lên ở thùng $A$}
\loigiai{ 
	\begin{itemchoice}
		\itemch  \textbf{Đúng}.\\
		 Ta có bảng số liệu ghép nhóm sau:
		\begin{center}
			\resizebox{0.8\linewidth}{!}{\begin{tabular}{|c|c|c|c|c|c|}
			 \hline
			Cân nặng (g) & $[250;260)$ & $[260;270)$ & $[270;280)$ & $[280;290)$  & $[290;300)$\\
			 \hline
			Giá trí đại diện & $255$ & $265$ & $275$  & $285$ & $295$\\
			 \hline
			Số táo ở thùng $A$  & $2$ & $4$ & $12$ & $4$ & $3$\\
			 \hline
			Số táo ở  thùng $B$ & $1$ & $3$ & $7$ & $10$ & $4$\\
			 \hline
		\end{tabular}}
	\end{center}
	Số trung bình của cân nặng các quả táo ở thùng $A$ là 
	\[\overline{x}_A=\dfrac{255\cdot2+265\cdot4+275\cdot12+285\cdot4+295\cdot3}{25}=275{,}8.\] 
	Số trung bình của cân nặng các quả táo ở  thùng $B$ là
	\[\overline{x}_B=\dfrac{255\cdot1+265\cdot3+275\cdot7+285\cdot10+295\cdot4}{25}=280{,}2.\] 
	Vậy số trung bình của cân nặng các quả táo ở thùng  $A$ nhỏ hơn số trung bình của cân nặng các quả táo ở thùng $B$.
	\itemch  \textbf{Đúng}.\\
	Gọi $M$ là~biến cố: \lq\lq Lấy được một  quả táo từ thùng $A$ có cân nặng từ $280$ g trở lên\rq\rq.\\
	Khi đó ta  có $\heva{&n(\Omega)=25\\&n(M)=4+3=7}\Rightarrow\mathrm{P}(M)=\dfrac{n(M)}{n(\Omega)}=\dfrac{7}{25}=0{,}28$.
	\itemch \textbf{Đúng}.\\ 
	Gọi $N$ là~biến cố: \lq\lq Lấy được một quả táo từ  thùng $A$ và một quả táo từ thùng $B$ sao cho hai quả táo lấy ra đều nặng từ $270$ g trở lên\rq\rq.\\
	Khi đó ta  có $\heva{&n(\Omega)=25\cdot25=625\\&n(N)=(12+4+3)\cdot(7+10+4)=399}\Rightarrow\mathrm{P}(M)=\dfrac{n(N)}{n(\Omega)}=\dfrac{399}{625}=0{,}6384$.
	 \itemch \textbf{Đúng}.\\
	Số  táo có cân nặng từ $280$ g trở lên ở thùng $B$ là~$10+4=14$.\\
	Số táo có cân nặng từ $280$ g trở lên ở thùng $A$  là
	$4+3=7$.\\ 
	Vậy số táo có cân nặng từ $280$ g trở lên ở thùng $B$ nhiều hơn số táo có cân nặng từ $280$ g trở lên ở thùng $A$. 
	\end{itemchoice}
 }
\end{ex}

% nguồn: Lop_12_Dai_so_Chuong_3
\begin{ex}%[2D3H1-4]%[ma:df88349c]
	Khi đo mắt cho học sinh khối 10 ở một trường THPT nhân viên y tế ghi nhận lại ở bảng  sau:
	\begin{center}\begin{tabular}{|c|c|c|c|c|c|}
		\hline 
		Thời  gian &$[0{,}25;0{,}75)$&$[0{,}75;1{,}25)$&$[1{,}25;1{,}75)$&$[1{,}75;2{,}25)$&$[2{,}25;2{,}75)$\\
		 \hline
		Số  lần&$25$&$32$&$14$&$12$&$4$\\
		 \hline
	\end{tabular}
\end{center}
 \choiceTF
	{\True Số trung bình của mẫu số liệu trên là  $1{,}14$}
	{\True Nhóm chứa mốt của số liệu là $\left[0{,}75;1{,}25\right)$} 
	{\True Mốt của mẫu số liệu là  $M_{\circ}=0{,}89$}
	{\True Trung vị của mẫu  số liệu là $M_e=1{,}039$}
 \loigiai{
	\begin{center}\begin{tabular}{|c|c|c|c|c|c|}
		 \hline
		Thời gian  &$[0{,}25;0{,}75)$&$[0{,}75;1{,}25)$&$[1{,}25;1{,}75)$&$[1{,}75;2{,}25)$&$[2{,}25;2{,}75)$\\
		 \hline
		Giá trị đại diện&$0{,}50$&$1{,}00$&$1{,}50$&$2{,}00$&$2{,}50$\\ 
		 \hline
		 Số lần&$25$&$32$&$14$&$12$&$4$\\
		 \hline
	\end{tabular}
	\end{center}
	\begin{itemchoice}
		\itemch  Số trung bình của mẫu số liệu trên là
		$$\dfrac{0{,}50\cdot25+1{,}00\cdot32+1{,}50\cdot14+2{,}00\cdot12+2{,}50\cdot4}{87}=1{,}14$$ 
		\itemch Nhóm chứa mốt  của số liệu là $[0{,}75;1{,}25)$.\\
		\itemch  Mốt của mẫu số liệu là $M_{\circ}=0{,}75+\dfrac{32-25}{(32-25)+(32-14)}(1{,}25-0{,}75)=0{,}89$.\\
		\itemch Gọi $x_1,x_2,\ldots{x_{87}}$ lần lượt là chỉ số mắt cận của các học sinh sắp xếp  theo thứ tự không giảm.\\
		Ta có $x_1,\ldots,x_{25}\in [0{,}25;0{,}75);x_{26},\ldots,x_{57}\in [0{,}75;1{,}25)$;  nên trung vị của mẫu là $x_{44}\in [0{,}75;1{,}25)$.\\
		Ta  xác định được $n=87,n_m=32,C=25,u_m=0{,}75;u_{m+1}=1{,}25$.\\
		 Nên khi đó $M_e=0{,}75+\dfrac{\dfrac{87}{2}-25}{32}\left(1{,}25-0{,}75\right)=1{,}039$.
	\end{itemchoice}
 }
\end{ex}

% nguồn: Lop_12_Dai_so_Chuong_3
\begin{ex}%[2D3V1-4]%[ma:7a0e10dc]
	Người ta đo độ dài thân của $80$ con cá rô được nuôi ở một khu vực. Kết quả được biểu diễn ở biểu đồ tần số tương  đối ghép nhóm sau (đơn vị cm).
	\begin{center}
		\resizebox{0.8\linewidth}{!}{\begin{tikzpicture}
		\draw[opacity=.25,thin,step=.2,cyan](0,0)  grid (7,3);
		\draw[opacity=1,cyan](0,0) grid (7,3); 
		\draw[stealth-stealth](0,3) node[left]{Tần số tương  đối (\%)}|-(7,0) node[below,anchor=west]{Độ dài (cm)};
		\foreach\x/\dientich[count=\i from 1] in  {0/27,0.25/28,1/29,2.5/30,0.75/31,0.5/32}{
			\draw[pattern=north west  lines,pattern color=cyan!65](\i-1,0) rectangle +(1,\x);
			\draw (\i,0)  node[below]{\dientich};
		} 
		\foreach \y [count=\i from 1] in {10,20,30,40,50}{ 
			 \draw (-1pt,\i/2)--(1pt,\i/2)
			(0,\i/2)  node[left]{$\y$};}
		\foreach \z [count=\i from  1] in {5,20,50,15,10}{
			 \draw (\i+0.5,\z/20) node[above] {$\z$};}
	\end{tikzpicture}
		} 
	\end{center}
	 \choiceTF
		{Số con cá có độ dài  thân nhỏ hơn $28$ cm là $5$}
		{\True Mốt của mẫu số liệu ghép nhóm thuộc khoảng $[29{,}4;29{,}5)$} 
		{\True Tứ phân vị thứ nhất của mẫu số liệu ghép nhóm là  $29$}
		{\True Khoảng tứ phân vị của mẫu số liệu ghép  nhóm là $1$}
	\loigiai{Ta có bảng tần số ghép nhóm 
		\begin{center}
			\begin{tabular}{|l|c|c|c|c|c|}
				\hline  Nhóm &{$[27;28)$}&{$[28;29)$}&{$[29;30)$}&{$[30;31)$}&{$[31;32)$}\\
				\hline Tần số & $4$ & $16$ & $40$ & $12$ & $8$  \\
				 \hline
			\end{tabular}
		\end{center}Cỡ mẫu $n=80$.\\
		\begin{itemchoice}
			 \itemch Sai. Số con cá có độ dài thân nhỏ hơn $28$ cm là $\dfrac{80\cdot5}{100}=4$.
			\itemch Đúng. Nhóm chứa mốt là $[29;30)$\\ 
			Mốt của mẫu số liệu ghép  nhóm là
			\[M_0=29+\dfrac{40-16}{(40-16)+(40-12)}\cdot(30-29)\approx29{,}46.\] 
			\itemch Đúng. Tứ phân vị thứ nhất của mẫu số liệu gốc là $\dfrac{1}{2}(x_{20}+x_{21})$. Mà $x_{20} \in[28;29)$,  $x_{21} \in[29;30)$.\\
			Do đó, tứ phân vị thứ nhất của mẫu số liệu ghép nhóm là  \[Q_1=29.\]
			\itemch 
			Đúng. Tứ phân vị thứ ba của mẫu số liệu gốc là $\dfrac{1}{2}(x_{60}+x_{61})$.\\  Mà $x_{60} \in[29;30)$, $x_{61} \in[30;31)$.\\ Do đó, tứ phân vị thứ ba của mẫu số liệu ghép nhóm là
			 \[Q_3=30.\]
			Vậy khoảng tứ  phân vị của mẫu số liệu ghép nhóm là
			 \[\Delta_Q=30-29=1.\]
	\end{itemchoice}}
\end{ex}

% nguồn: Lop_12_Dai_so_Chuong_3
\begin{ex}%[2D3V1-4]%[ma:36bfd6d5]
	Bảng  sau đây biểu diễn mẫu số liệu nhóm về đường kính của một số quả cam sành Hà Giang được lựa chọn ngẫu nhiên từ một lô hàng (đơn vị mm).
	\begin{center}
		\begin{tabular}{|l|c|c|c|c|c|}
			 \hline Nhóm &{$[75;77)$}&{$[77;79)$}&{$[79;81)$}&{$[81;83)$}&{$[83;85)$}\\
			\hline Tần số &  $12$ & $25$ & $38$ & $20$ & $5$ \\
			 \hline
		\end{tabular}
	\end{center}\choiceTF
	{\True Số phần tử của mẫu  (cỡ mẫu) là $n=100$}
	{Khoảng biến thiên của mẫu số liệu ghép nhóm trên  là $8$}
	{Tứ phân vị thứ ba của  mẫu số liệu ghép nhóm trên là ${{Q}_{3}}=83$}
	{\True Khoảng tứ phân vị  của mẫu số liệu ghép nhóm trên là ${{\Delta }_{Q}}=2,96$}
	\loigiai{\begin{itemchoice}
		\itemch Đúng. Cỡ  mẫu $n=100$.
		\itemch Sai. Khoảng biến thiên của mẫu  số liệu ghép nhóm trên là $85-75=10$.
		\itemch Sai. Gọi $x_1; x_2; \ldots; x_{100}$ là mẫu số liệu gốc về đường kính của một số  quả cam sành Hà Giang được lựa chọn ngẫu nhiên từ một lô hàng được xếp theo thứ tự không giảm.\\
		Ta  có
		\begin{listEX}[3]
			\item $x_1, \ldots,  x_{12} \in[75;77)$
			\item $x_{13}, \ldots, x_{37}  \in[77;79)$
			 \item $x_{38}, \ldots, x_{75} \in[79;81)$
			\item  $x_{76}, \ldots, x_{95} \in[81;83)$
			\item  $x_{96}, \ldots, x_{100} \in[83;85)$.
		\end{listEX}
		 Tứ phân vị thứ ba của mẫu số liệu gốc là $\dfrac{1}{2}(x_{75}+x_{76})$.\\ Mà: $x_{75} \in[79;81)$, $x_{76} \in[81;83)$.\\ Do đó, tứ phân vị thứ ba của mẫu số liệu ghép nhóm là
		\[Q_3=81.\] 
		\itemch Đúng. Tứ phân vị thứ nhất của  mẫu số liệu gốc là $\dfrac{1}{2}(x_{25}+x_{26}) \in[77;79)$.\\
		Do  đó, tứ phân vị thứ nhất của mẫu số liệu ghép nhóm là
		\[Q_1=77+\dfrac{\dfrac{100}{4}-12}{25}  \cdot(79-77)=78{,}04.\]\\
		Khoảng tứ phân vị của mẫu số liệu  ghép nhóm là
		 \[\Delta_Q=81-78{,}04 =2{,}96.\]
	\end{itemchoice}
} 
\end{ex}

% nguồn: Lop_12_Dai_so_Chuong_3
\begin{ex}%[2D3V1-4]%[ma:5135bc0f]
	Bảng sau đây biểu diễn lượng mưa trung bình đo được tại một trạm  quan trắc đặt tại Nam Định trong các năm từ $2007$ đến $2023$ (đơn vị mm).
	\begin{center}
		\resizebox{0.8\linewidth}{!}{\begin{tabular}{|l|c|c|c|c|c|c|c|c|}
		\hline  {$1\,114$} &{$1\,087$}&{$1\,800$}&{$1\,643{,}6$}&{$1\,461{,}4$}&,{$1\,767{,}2$}&{$1\,772{,}8$}&{$1\,757{,}3$}&{1\,721{,}4}\\
		\hline  {$1\,349{,}7$}&{$1\,612{,}3$}&{$2\,318{,}3$}&{$1\,800{,}1$}&{$1\,265$}&{$1\,641{,}5$}&{$2\,227{,}3$}&{$2\,542{,}4$}&{}\\
		\hline 
	\end{tabular}}
\end{center}
Người ta lập bảng dữ liệu ghép nhóm cho mẫu số  liệu trên. Bảng gồm 5 nhóm có độ dài bằng nhau, nhóm đầu tiên là $[1\,050;1\,350)$. Sử dụng bảng số liệu đó, ta có
\choiceTF 
	{\True  Đầu mút trái của nhóm cuối cùng là $2\,250$}
	 {Khoảng biến thiên của mẫu số liệu ghép nhóm là $1\,050$}
	{\True Tứ phân vị thứ nhất của mẫu số liệu ghép nhóm  nhỏ hơn $1\,452$}
	{Khoảng tứ phân vị của mẫu số liệu  ghép nhóm lớn hơn $519$}
\loigiai{ 
	Bảng  tần số ghép nhóm
	\begin{center}
		\resizebox{0.8\linewidth}{!}{\begin{tabular}{|l|c|c|c|c|c|}
		\hline Nhóm &{$[1\,050;1\,350)$}&{$[1\,350;1\,650)$}&{$[1\,650;1\,950)$}&{$[1\,950;2\,250)$}&{$[2\,250;2\,550)$}\\ 
		\hline Gía trị đại diện &{$1\,200$}&{$1\,500$}&{$1\,800$}&{$2\,100$}&{$2\,400$}\\ 
		\hline Tần số &  $4$ & $4$ & $6$ & $1$ & $2$ \\
		\hline 
	\end{tabular}}
	\end{center}
	\begin{itemchoice}
		\itemch Đúng. Đầu mút  trái của nhóm cuối cùng là $2\,250$.
		\itemch Sai. Khoảng biến  thiên của mẫu số liệu ghép nhóm là $2550-1050=1500$.
		\itemch Đúng. Tứ phân vị thứ  nhất của mẫu số liệu ghép nhóm là \[Q_1=1368{,}75<1\,452.\]
		\itemch Sai. Tứ phân vị thứ ba của mẫu số  liệu ghép nhóm là
		\[Q_3=1887{,}5.\] 
		Vậy khoảng tứ phân vị  của mẫu số liệu ghép nhóm là
		\[\Delta_Q=518{,}75<519.\] 
\end{itemchoice} }
\end{ex}

% nguồn: Lop_12_Dai_so_Chuong_3
\begin{ex}%[2D3V1-4]%[ma:0f0b3d31]
	Time on Site hay còn được gọi là  thời gian trên trang là thời gian khách hàng lưu lại trên một website. Bảng sau đây ghi lại thời gian trên trang của một số khách hàng ở một website (đơn vị giây)
	\begin{center}
		\begin{tabular}{|l|c|c|c|c|c|}
			\hline Nhóm &{$[0;60)$}&{$[60;120)$}&{$[120;180)$}&{$[180;240)$}&{$[240;300)$}\\ 
			\hline Tần số & $125$ &  $100$ & $35$ & $30$ & $10$ \\
			 \hline
		\end{tabular}
	\end{center}
	\choiceTF 
		{\True Khoảng biến thiên của mẫu số  liệu ghép nhóm trên là $300$}
		{Tứ phân vị thứ nhất và tứ phân vị  thứ hai của mẫu số liệu ghép nhóm trên thuộc cùng một nhóm}
		{\True Tứ phân vị thứ ba của mẫu số liệu ghép nhóm trên thuộc nhóm  $[120;180)$}
		{Khoảng tứ phân vị của mẫu số liệu ghép nhóm  trên là $144$}
	\loigiai{\begin{itemchoice}
		\itemch  Đúng. Khoảng biến thiên $300-0=300$.
		 \itemch Sai. Cỡ mẫu $n=300$.\\
		Gọi $x_1; x_2; \ldots; x_{300}$ là mẫu số  liệu gốc về thời gian trên trang của một số khách hàng ở một website được xếp theo thứ tự không giảm.\\
		Ta  có
		\begin{listEX}[3]
			\item $x_1, \ldots,  x_{125} \in[0;60)$
			 \item $x_{126}, \ldots, x_{225} \in[60;120)$
			\item  $x_{226}, \ldots, x_{260} \in[120;180)$
			\item $x_{261},  \ldots, x_{290} \in[180;240)$
			\item $x_{291}, \ldots,  x_{300} \in[240;300)$.
		\end{listEX}
		Tứ phân vị  thứ nhất của mẫu số liệu gốc là $\dfrac{1}{2}(x_{75}+x_{76}) \in[0;60)$.\\
		Do  đó, tứ phân vị thứ nhất của mẫu số liệu ghép nhóm là
		\[Q_1=0+\dfrac{\dfrac{100}{4}-0}{125}  \cdot(60-0)=12.\]\\
		Tứ phân vị thứ  hai của mẫu số liệu gốc là $\dfrac{1}{2}(x_{150}+x_{151}) \in[60;120)$.\\
		Tứ phân vị thứ nhất và tứ phân  vị thứ hai của mẫu số liệu ghép nhóm trên không thuộc cùng một nhóm.
		\itemch Đúng. Tứ phân vị thứ ba  của mẫu số liệu gốc là $\dfrac{1}{2}(x_{225}+x_{226})$.\\ Mà: $x_{225} \in[60;120)$, $x_{226} \in[120;180)$.\\ Do đó, tứ phân vị thứ ba của mẫu số liệu ghép nhóm là
		 \[Q_3=120.\]
		\itemch Sai. Khoảng tứ phân vị của mẫu số liệu  ghép nhóm là
		\[\Delta_Q=Q_3-Q_1=120-12  =108.\]
\end{itemchoice}}
\end{ex}

% nguồn: Lop_12_Dai_so_Chuong_3
\begin{ex}%[2D3V1-4]%[ma:e435378b]
	Bạn An và bạn Bình làm thí nghiệm trồng cây. Mỗi bạn trồng $ 40$ cây cần tây trong cốc, phần gốc của các cây khi bắt đầu trồng  đều dài $4$ cm. \textit{Bảng 13} và \textit{Bảng 14} lần lượt biểu diễn mẫu số liệu ghép nhóm về số liệu thống kê chiều cao của các cây (đơn vị: centimét) mà bạn An và bạn Bình trồng sau $5$ tuần.
	\begin{center}
		\begin{tabular}{|c|c|}
			\hline 
			Nhóm &  Tần số \\
			\hline 
			$[20; 25)$ &  $2$ \\
			$[25; 30)$  & $16$ \\
			 $[30; 35)$ & $20$ \\
			$[35; 40)$ &  $2$ \\
			 \hline
			\multicolumn{2}{|c|}{\textit{Bảng 13}} \\ 
			\hline 
		\end{tabular}
		\begin{tabular}{|c|c|}
			\hline 
			Nhóm  & Tần số \\
			\hline 
			$[20;  25)$ & $5$ \\
			$[25; 30)$ &  $9$ \\
			$[30;  35)$ & $25$ \\
			 $[35; 40)$ & $1$ \\
			 \hline
			\multicolumn{2}{|c|}{\textit{Bảng  14}} \\
			 \hline
		\end{tabular}
	\end{center}
	\choiceTF 
		{Chiều cao trung bình của mỗi cây  do hai bạn An và Bình trồng không bằng nhau}
		{\True Khoảng biến thiên của cả hai mẫu  số liệu trên là $20$}
		{\True Khoảng tứ  phân vị của mẫu số liệu ở \textit{Bảng 13} là $5{,}5$}
		{Chiều cao của các cây mà  bạn Bình trồng đồng đều hơn các cây mà bạn An trồng}
	 \loigiai{
		Ta có  bảng biểu diễn mẫu số liệu ghép nhóm về số liệu thống kê chiều cao của các cây mà bạn An trồng là
		\begin{center}
			\begin{tabular}{|c|c|c|c|}
				\hline 
				Nhóm & Giá trị đại diện  & Tần số & Tần số tích lũy\\
				\hline 
				 $[20; 25)$ & $22{,} 5$& $2$ & $2$ \\
				$[25; 30)$ & $27{,} 5$  & $16$ & $18$\\
				$[30; 35)$ & $32{,} 5$ & $20$ & $38$\\ 
				$[35; 40)$ & $37{,} 5$& $2$ & $40$\\ 
				\hline 
			\end{tabular}
		\end{center}
		Suy ra chiều cao trung bình của mỗi cây mà  An trồng là \[\overline{x}_A= \dfrac{2 \cdot 22{,}5 + 16 \cdot 27{,}5 + 20 \cdot 32{,}5 + 2 \cdot 37{,} 5}{40}= 30{,} 25 \, (\text{cm}).\]
		Bảng biểu diễn mẫu số liệu ghép nhóm về số liệu thống kê chiều  cao của các cây mà bạn Bình trồng là
		\begin{center}
			\begin{tabular}{|c|c|c|c|}
				 \hline
				Nhóm &  Giá trị đại diện & Tần số & Tần số tích lũy\\
				 \hline
				$[20; 25)$ & $22{,} 5$& $5$ &  $5$ \\
				$[25;  30)$ & $27{,} 5$ & $9$ & $14$\\
				$[30; 35)$ & $32{,} 5$  & $25$ & $39$\\
				$[35; 40)$ & $37{,}  5$& $1$ & $40$\\
				\hline 
			\end{tabular}
		\end{center}
		Suy ra chiều cao  trung bình của mỗi cây mà Bình trồng là \[\overline{x}_B= \dfrac{5 \cdot 22{,}5 + 9 \cdot 27{,}5 + 25 \cdot 32{,}5 + 1\cdot 37{,} 5}{40}= 30{,} 25 \, (\text{cm}).\]
		\begin{itemchoice}
			\itemch {\bf  Sai}. \\
			Vì $\overline{x}_A=  \overline{x}_B= 30{,} 25$ cm.
			\itemch {\bf Đúng}.  \\
			Vì khoảng biến thiên của cả hai mẫu số liệu là  $40-20=20$.
			\itemch {\bf Đúng}.\\ 
			Xét mẫu số liệu ở \textit{Bảng  13}.\\
			Ta có $\dfrac{n}{4}  = 10$.\\
			Suy ra nhóm đầu tiên có tần số  tích lũy lớn hơn hay bằng $10$ là $ [25; 30)$.\\
			Suy ra tứ phân vị  thứ nhất của mẫu số liệu này là \[Q_1= 25 +\dfrac{10 - 2}{16} \cdot 5=27{,}5.\]
			Ta có $\dfrac{3n}{4}= 30$.\\ 
			Nhóm đầu tiên có tần số tích lũy lớn hơn  hay bằng $30$ là nhóm $[30; 35)$.\\
			Suy ra tứ phân vị thứ ba  của mẫu số liệu này là \[Q_3= 30 +\dfrac{30- 18}{20} \cdot 5=33.\]
			Vậy khoảng tứ phân vị của mẫu số liệu này là $Q_3-Q_1= 33- 27{,}5 =  5{,}5 $.
			\itemch {\bf Sai}. \\ 
			Ta có \[s^2_A = \dfrac{2 \cdot 22{,}5^2 + 16 \cdot 27{,}5^2 + 20 \cdot 32{,}5^2 + 2 \cdot 37{,} 5^2}{40}  - (30{,} 25)^2= 11{,} 1875.\]
			\[s^2_B=\dfrac{5 \cdot 22{,}5^2 + 9 \cdot 27{,}5^2 + 25 \cdot 32{,}5^2 + 1\cdot 37{,} 5^2}{40}-(30{,}25)^2=  13{,}6875.\]
			Vì $s_A^2 <s_B^2$ nên chiều cao mà các cây bạn An trồng đồng đều hơn các cây  mà bạn Bình trồng.
	\end{itemchoice}}
\end{ex}


%%%%%%%%%%%%%%%% TRẢ LỜI NGẮN %%%%%%%%%%%%%%%%

% nguồn: Lop_12_Dai_so_Chuong_3
\begin{bt}%[2D3H1-4]%[ma:6fcbb8b3]
	 Bác Tuấn kiểm tra cân nặng của 40 quả trứng được lựa chọn ngẫu nhiên từ một trang trại và ghi kết quả vào bảng dữ liệu ghép nhóm sau:
	\begin{center}
		\begin{tabular}{|c|c|c|c|c|c|}
			\hline 
			Cân nặng (gam) & $[75;80)$ & $[80;85)$ & $[85;90)$ & $[90;95)$  & $[95;100)$ \\ \hline
			Tần số tương đối & $25\%  $ & $35\%$ & $25\% $ & $10\% $ & $5\%$ \\ \hline
		\end{tabular}
	\end{center}
	Tính tỉ số của khoảng tứ phân vị và khoảng biến thiên của  mẫu số liệu trên.
	\shortans{0,32} 
	 \loigiai{
		Bảng tần số ghép  nhóm
		\begin{center}
			\begin{tabular}{|c|c|c|c|c|c|}
				\hline 
				Nhóm & $[75;80)$ & $[80;85)$ & $[85;90)$ & $[90;95)$ & $[95;100)$ \\  \hline
				 Tần số & $10$ & $14$ & $10$ & $4$ & $2$ \\ \hline
			\end{tabular}
		\end{center}
		Khoảng biến thiên của mẫu số liệu ghép nhóm là $100 -75 =  25$ (g).\\
		Tứ phân vị thứ nhất của mẫu số liệu ghép nhóm là $Q_1=80$; tứ  phân vị thứ ba của mẫu số liệu ghép nhóm là $Q_3=88$.\\
		Khoảng tứ phân vị của mẫu số liệu ghép  nhóm là $88-80=8$.\\
		Tỉ  số của khoảng tứ phân vị và khoảng biến thiên là $8 \div 25 = 0{,}32$.
	} 
\end{bt}
